\section{Related literature}
\label{sec:lit}

The literature on university technology transfer is now large, covering patenting and licensing, academic entrepreneurship, university--industry engagement, spin-offs, institutional missions, and the organizations that support commercialization \citep{CunninghamMenterStarke2025,PerkmannEtAl2013,PerkmannEtAl2021}. Universities also rely on a wider infrastructure that includes TTOs, incubators, science parks, and university venture funds \citep{GoodEtAl2019}. Our focus is narrower: the formal IP system through which inventions enter the university, are evaluated and protected, and eventually reach outside users through licensing.

Much of the early work on that system focused on ownership and incentives. The Bayh--Dole Act allowed universities and other federal research contractors to elect title to qualifying federally funded inventions, and a large literature followed the growth of university patenting and licensing that came with institutional management of research inventions \citep{MoweryEtAl2001,HendersonJaffeTrajtenberg1998,Sampat2006,GrimaldiEtAl2011}. Changes in ownership rules can alter researcher incentives and behavior. When Norway ended the ``professor's privilege'' and shifted ownership from researchers to universities, patenting and entrepreneurship by university researchers fell \citep{HvideJones2018}. Related work documents the broader European move toward institutional ownership and shows how different ownership arrangements change the incentives facing academic inventors \citep{GeunaRossi2011,KenneyPatton2009}.

Researchers have also asked whether universities can influence commercialization through more direct financial incentives. \citet{LachSchankerman2008} linked larger inventor royalty shares to higher licensing income, while \citet{DiGregorioShane2003} and \citet{ArqueCastellsEtAl2016} studied equity, royalty sharing, startup formation, and inventive effort. But the evidence here is less settled than it once appeared. Using corrected and expanded U.S. policy data, \citet{OuelletteTutt2020} show that the earlier relationship between royalty shares and licensing income was driven by coding errors and find no compelling evidence that larger royalty shares significantly affected university patenting or commercialization outcomes. The result is especially relevant here because even a seemingly direct policy lever can be difficult to isolate when universities differ in research portfolios, institutional capacity, governance, and commercialization strategy.

A different strand of the literature looks inside the technology-transfer process itself. Studies of TTOs emphasize differences in staffing, incentive structures, organizational routines, and relationships with industry \citep{SiegelWaldmanLink2003,SiegelWaldmanAtwater2004,ThursbyKemp2002,MarkmanEtAl2005,OsheaEtAl2005}. \citet{JensenThursbyThursby2003} model the TTO as an intermediary between faculty and the university administration, with objectives of its own that can affect both disclosure and licensing. \citet{DebackereVeugelers2005} emphasize the governance structures and decision processes through which TTOs connect university research to firms. Institutional objectives also matter: public universities with local-development goals have been shown to license differently from private universities \citep{BelenzonSchankerman2009}. More recent work places the TTO within the broader task of intellectual-property management and argues that its role includes helping move university knowledge into use, not only securing patents and capturing revenue \citep{HolgerssonEtAl2019,AaboenHolgersson2019}.

This organizational work also directs attention to the relationship between researchers and the TTO. Many potentially commercial inventions never enter the formal system \citep{JensenThursbyThursby2003}. Faculty disclosure depends on local norms, what colleagues do, and how researchers view the technology-transfer office \citep{OwenSmithPowell2001,BercovitzFeldman2008}, and part of the historical growth in university licensing has been linked to greater faculty willingness to disclose \citep{ThursbyThursby2002}. TTOs themselves must establish legitimacy with academics as well as university management \citep{OKaneEtAl2015}. Recent work also shows that faculty sometimes bypass formal technology-transfer channels, with previous experiences with the office, peer views, their own commercialization expertise, and outside partners all playing a role \citep{HalilemDiop2025}. Perceptions of fairness matter as well: in a survey of 1,329 academic entrepreneurs at 25 U.S. research universities, \citet{WaldmanEtAl2022} find that organizational-justice perceptions are positively associated with intentions to use formal technology-transfer channels and negatively associated with intentions to use informal, noncompliant channels. \citet{Siegel2026} points to these fairness perceptions and to bypassing as part of a broader agenda of micro-organizational and behavioral research on technology transfer. Together, this work moves the discussion beyond financial incentives toward the institutional relationship through which researchers engage with commercialization.

That relationship continues after an invention is disclosed. University inventions are often licensed while still at an early stage of development, and successful commercialization can require continued inventor involvement after the formal disclosure and licensing process has begun \citep{JensenThursby2001}. It is therefore useful to distinguish two broad margins. The first is an entry margin: whether an inventor brings an invention into the university's formal technology-transfer system. The second is a post-disclosure margin: what happens once the invention is inside that system, as the university evaluates it, decides whether to seek protection, searches for potential licensees, negotiates with firms, and continues to work with the inventor.

Our paper builds on these literatures. Much of the existing work measures ownership arrangements, financial incentives, TTO resources, organizational practices, or researcher perceptions. Less attention has been given to how universities formally communicate the governance relationship between the institution and its inventors. An IP policy does more than specify who owns an invention or how royalties are divided. It also states obligations, allocates authority, describes procedures, and tells inventors what forms of assistance the institution provides.

The companion paper \citep{SharmaPCSI2026} measures this dimension with the Policy Communication Stance Index (PCSI), using text-as-data methods \citep{LoughranMcDonald2011,GentzkowKellyTaddy2019,GrimmerRobertsStewart2022}. PCSI captures how rights, obligations, institutional authority, and assistance are presented in formal IP-policy text. It does not measure whether a policy is legally more generous, whether faculty perceive the university as supportive, or whether the TTO actually behaves in the way the policy describes. Its purpose is narrower: to measure the communicative stance of formal IP governance consistently across universities and over time.

Why might written policy be related to technology-transfer outcomes? One possibility runs through researchers themselves. Work on administrative burden shows that formal procedures can impose learning, compliance, and psychological costs on the people who navigate them \citep{MoynihanHerdHarvey2015,HerdMoynihan2018}. In this setting, policy may help shape how researchers understand what they must disclose, what assistance is available, and how they are expected to work with the university. Policy may also matter through the TTO. Formal IP rules define responsibilities, authority, and procedures within the same organizational setting in which staff evaluate inventions and interact with inventors. The literature on TTO organization gives us good reason to expect those routines and relationships to matter \citep{JensenThursbyThursby2003,DebackereVeugelers2005,HolgerssonEtAl2019}. This channel does not require every inventor to read the policy directly. Formal policy provides part of the governance framework within which TTO procedures and inventor interactions take place.

There is an important alternative interpretation. Written policy may not itself be the mechanism. Institutional theory has long emphasized that formal rules can be partly decoupled from organizational practice \citep{MeyerRowan1977}. A rewrite may therefore be one visible part of a larger change in leadership, TTO organization, workflows, or commercialization strategy. A change in PCSI may then mark a revision episode without being the cause of the outcomes that follow.

These possibilities are not mutually exclusive, and our data cannot separate them. They do, however, help organize the empirical analysis. If the main change is at the point where researchers decide whether to enter the formal technology-transfer system, invention disclosures are the natural early outcome to examine. If the relevant changes occur later in the process, differences may instead appear in patent applications, patents issued, or licensing. Broader organizational reforms could affect several stages at once.

This is why we examine several outcomes rather than a single measure of technology-transfer performance. Invention disclosures capture entry into the formal system. New patent applications and the application-to-disclosure filing margin capture movement from disclosure toward patent protection, although the institution-year filing margin is not an invention-level conversion probability \citep{ChoudhryPonzio2020,Hall2021}. Patents issued lie farther downstream. Licenses and options executed capture the stage at which university technologies are formally transferred to outside users. Looking across these outcomes allows us to ask where a post-revision pattern is strongest, even when the aggregate data cannot tell us which mechanism produced it.

Most university IP-policy changes are much smaller than the large ownership reforms studied in the quasi-experimental literature. Universities periodically rewrite their policies while leaving institutional ownership in place. These revisions may change disclosure procedures, revenue-sharing rules, conflict-of-interest provisions, equity arrangements, assignment language, appeals, or the way these rules are presented to inventors. Yet ordinary within-university revisions have received little attention as dated events around which commercialization outcomes can be studied.

We use PCSI to classify the direction of these revisions and follow technology-transfer outcomes around them. This shifts the comparison from differences in individual policy provisions across universities toward policy change within institutions. Our goal is to ask where outcomes diverge around revision episodes, not to estimate the causal effect of particular words, clauses, or communicative styles.
